\section{Related Work}
\label{sec:Related}

\begin{table}[t]
\centering
\caption{Comparison of zkPHIRE with Prior ZKP Accelerators.\\N = NTT, S = SumCheck, M = MSM. Areas are in 7nm.}
\label{tab:comp_table}
\rowcolors{2}{rowgray}{white}
\resizebox{\columnwidth}{!}{%
\begin{tabular}{>{\bfseries}l|c c c c}
\hline
\textbf{Metric} & \textbf{NoCap} & \textbf{SZKP+} & \textbf{zkSpeed+} & \textbf{zkPHIRE} \\
\hline
Workload & Scaled-up AES & \begin{tabular}[c]{@{}c@{}}Rollup 25\\Pvt Tx\end{tabular} & \begin{tabular}[c]{@{}c@{}}Rollup 25\\Pvt Tx\end{tabular} & \begin{tabular}[c]{@{}c@{}}Rollup 25\\Pvt Tx\end{tabular}
\\
Protocol & Spartan+Orion & Groth16 & HyperPlonk & HyperPlonk \\
Main Kernels & N \& S & N \& M & S \& M & S \& M
\\
Gates & $2^{24}$ & $2^{24}$ & $2^{24}$ & $2^{19}$
\\
Encoding & R1CS & R1CS & \begin{tabular}[c]{@{}c@{}}Plonk\\ (Vanilla)\end{tabular} & \begin{tabular}[c]{@{}c@{}}Plonk\\ (Jellyfish)\end{tabular} \\
Proof Size & 8.1 MB & 0.18 KB & 5.09 KB & 4.41 KB \\
Setup & none & circuit-specific & universal & universal \\
Prime & fixed & arbitrary & arbitrary & fixed \\
Bitwidth & 64 & 255/381 & 255/381 & 255/381 \\
SW Prover (s) & 94.2 & 51.18 & 145.5 & 6.161 \\
HW Prover (ms) & 151.3 & 28.43 & 151.973 & 3.874 \\
SW Verifier (ms) & 134 & 4.2 & 26 & 19 \\
Chip Area (mm\textsuperscript{2}) & 38.73 & 353.2 & 366.46 & 294.32 \\
\# Modmuls & 2432 & 1720 & 1206 & 2267 \\
Modmul (mm\textsuperscript{2}) & 0.002 & 0.133 / 0.314 & 0.133 / 0.314 & 0.073/0.162 \\
Power (W) & 62 & $>$220 & 171 & 202 \\
\hline
\end{tabular}
}
\end{table}


Prior work on cryptographic acceleration has largely focused on Multi-Party Computation and Homomorphic Encryption \cite{kim2022_BTS, sharp, rpu, haac, karthik1, karthik2, ciflow, mo2023towards, ark, jha2024deepreshape, osiris, anaheim}.
Research on ZKP acceleration is relatively newer, and has focused primarily on MSMs and NTTs~\cite{priormsm, distmsm, cuzk, rezk, legozk, unizk, pipezk, gzkp, gypso, MSMAC, tches_ntt_msm, zhang2025code, zhang2025micro, elastic_msm, sam, myosotis, graz, prophet}.
The community has broadened its focus to include other ZKP kernels such as the SumCheck protocol and Merkle tree-based commitments.
Prior work has explored accelerating these kernels on GPUs~\cite{batchzk} and ASICs~\cite{nocap, zkspeed2025}.
We compare zkPHIRE with three ASICs that accelerate proof-generation end-to-end: SZKP, NoCap, and zkSpeed.

\textbf{Protocol Acceleration}:
\textbf{SZKP} is a Groth16~\cite{groth16} accelerator that speeds up all MSMs, including sparse G2 MSMs, achieving a geomean speedup of 493× over CPU. While a direct comparison is not apples-to-apples due to differing proof sizes and trusted setup assumptions, Groth16 and HyperPlonk target similar application domains requiring small, efficiently verifiable proofs. Unlike Groth16, HyperPlonk supports a universal setup and avoids an expensive trusted setup ceremony~\cite{ceremony}. Following zkSpeed, we report results for a scaled SZKP+ configuration that adopts more efficient MSMs and the BLS12-381 curve.
\textbf{NoCap} accelerates Spartan+Orion proofs using a vector architecture. It relies on Merkle tree-based commitments, resulting in larger proofs and higher verifier costs, and is better suited for settings where proof succinctness is less critical or number of verifiers is small.
\textbf{zkSpeed} is, to our knowledge, the only prior accelerator targeting HyperPlonk end-to-end, achieving a geometric mean speedup of 801× over CPU. For our comparison here, we use zkSpeed+ as we have done earlier. 

\autoref{tab:comp_table} summarizes the comparison across the four accelerators and reports iso-workload performance numbers.
We use numbers for a scaled-up AES workload from NoCap's evaluation as it has a similar number of gates to the Rollup 25 workload. 
We use the same area scaling factors used by zkSpeed to scale our area and NoCap's area to 7nm.
In our experiments, we assume a 1:1 mapping between R1CS constraints and Vanilla constraints. If Vanilla circuits require more constraints, zkPHIRE’s Jellyfish gains relative to Vanilla gates would be larger, scaling roughly linearly at iso-workload.
zkPHIRE's proving time is 39$\times$, 7$\times$, and 39$\times$ faster than NoCap, SZKP+, and zkSpeed+, respectively.
Jellyfish's advantages are application-dependent: not all workloads map efficiently to Jellyfish gates. In such cases, different custom gates (which zkPHIRE can support) may yield faster proving times. 
Finally, while SZKP+ still yields the smallest proofs, it suffers from a major usability drawback due to the costly circuit-specific trusted setup.

\textbf{Limitations of prior work}:
Prior work \cite{nocap} accelerates SumCheck using vector processors, but vector-style reductions introduce significant overhead: reductions over length-$M$ vectors require $\log_2(M)$  serialized folding steps, repeated register-file accesses, and additional reductions for higher-degree extensions. zkPHIRE avoids these bottlenecks by using fused, tree-structured reduction pipelines with minimal data movement.
While prior work \cite{nocap, legozk} can scale to workloads of $2^{30}$ constraints, their reliance on NTT-based matrix encodings introduces substantial latency overhead, e.g., off-chip transposes and recursive decompositions.
In contrast, zkPHIRE has straightforward memory accesses, and given the streaming architecture, can scale well beyond $2^{30}$ constraints, as long as the working set fits in off-chip memory. For ZKML or other problems of this scale, zkPHIRE would need better data orchestration schemes (e.g. fetching from the cloud). We leave this for future work.

\section{Conclusion}
This paper presents zkPHIRE, the first accelerator to support programmable SumCheck for succinct proof generation at very large scales, enabling efficient ZKPs over custom, high-degree gates.
On HyperPlonk baselines, zkPHIRE achieves iso-application geomean speedups of 
$11.87\times$ over prior ASICs and $1486\times$ over CPU.
Its design abstracts SumCheck into common primitives—updates, extensions, and products—and exploits parallelism within tables, across tables, and across evaluation products to scale across diverse polynomial structures.
By streaming polynomial tables through compact scratchpads and coupling compute units via a Multifunction Forest, zkPHIRE improves resource sharing and reduces SRAM and datapath overheads. Finally, a systematic design-space exploration balances parallelism, area efficiency, and performance to support a wide range of polynomials and custom gate constructions.
